
Reinforcement learning from human feedback (RLHF) aligns language models by training on human preference judgments (Ouyang et al., 2022; Bai et al., 2022). Current evaluation focuses on AI-side metrics: response quality (AlpacaEval, MT-Bench), factual accuracy (TruthfulQA), and safety properties. The bidirectional alignment framework \citep{shen2024bidirectional} identifies a complementary dimension: $Human\rightarrow AI$ alignment, which concerns whether alignment processes preserve human agency, critical evaluation capacity, and decision autonomy. No computational metrics currently operationalize this dimension.

Linguistic research documents associations between specific textual features and psychological constructs in human language. Juanchich et al. (2017) found that modal verbs (could, might, should) correlate with perceived autonomy attribution in controlled experiments. Biber et al. (1999) documented hedging phrases (appears, seems, possibly) as markers of epistemic stance. These findings establish that linguistic features can indicate agency-related constructs in human-authored text. Whether these associations transfer to AI-generated responses in RLHF contexts remains untested.

This study addresses the question: Can linguistic agency markers validated in human psychology serve as computational proxies for measuring agency preservation in RLHF evaluation? We employ a multi-criterion validation protocol testing (1) extraction feasibility, (2) effect size magnitude, (3) internal consistency across markers, and (4) cross-dataset replication. The validation approach treats proxy transfer as an empirical hypothesis rather than a theoretical assumption.
